% Section 1: Introduction and programme
\label{sec:intro-body}

\subsection{The question this paper asks}

Papers~0--IV of this series follow one expression through four successively coarser
stages: an expression tree, a marked history, the operator that the history induces,
and the endpoint value where it lands \cite{Yuan2026AEGPaperZero,Yuan2026AEGFoundations}.
Paper~IV then asks when a quotient of that data is still sufficient for future
continuation, and separates a semantic residual from an operational trace
\cite{Yuan2026AEGCondensation}.

This paper asks a different question, which the earlier papers deliberately do not
answer.  Suppose one and the same carrier is to be read in three coordinated ways ---
say as a process, as a space, and as a construction --- and suppose each reading is
a typed interpretation with its own direction of information flow.  What is the
minimal structure that lets such readings be compared at all?  What is necessarily
lost in the comparison, and how is that loss to be accounted for?  And what is the
syntax in which the places of loss --- the places where a reading is *not yet*
determined --- can be written down, filled, and checked?

The answer developed here has two halves, and they are the two parts of the paper.

\subsection{Part I in one paragraph}

Part~I develops a theory of \emph{triadic characteristics}.  Its primitive is not a
value but an \emph{observer contract}: a finite datum fixing what is seen, when two
things are declared indistinguishable, which projections are available, and what
budget is granted.  A characteristic of a process is then an operational quotient of
that process together with a residual, and the residual is not a bookkeeping
afterthought: it is the part of the construction that the quotient cannot see.

Given that framing, several classical-looking questions become sharp.

\begin{itemize}
\item Among the one-parameter quadratic scalar families
  $A_\kappa=\R[j]/(j^2-\kappa)$, which one does the geometry force?  The answer is
  that nothing in rank-one characteristic recovery, conjugation, or a multiplicative
  quadratic norm forces any of them; what selects $\C$ is a global negative-inverse
  involution on the punctured carrier, and this is equivalent, under one explicit
  bridge hypothesis, to the existence of an oriented order-four exchange lift
  $J$ satisfying $J^2=-I$ whose projective class satisfies $[J]^2=[I]$.
  \emph{Chirality} is precisely the central sign that projectivization forgets.
\item Where does the six-state carrier $V\oplus\Lambda^2 V$ come from?  Not from a
  postulate: on a declared finite cut interval, cut-membership cochains together
  with oriented interchange generate the Clifford anticommutator, whereas
  interchange alone provably cannot.
\item What survives when one finite characteristic algebra is read temporally,
  spatially, and constructively?  A stable variance signature, inference as
  constrained intersection of observation fibres, and a hard separation between four
  numeric invariants that a naive theory conflates.
\item What is the right notion of \emph{crossing}?  Not a transport arrow, but a
  cospan with a resolution torsor; monodromy is then the difference of two
  resolutions, and the Legendre monodromy is \emph{derived} from the quadratic base
  change rather than inserted by hand.
\end{itemize}

Part~I also contains the paper's principal \emph{negative} calibration.  On the
finite Legendre prism the total defect decomposes into a central sign channel and a
history-suffix channel; therefore no irreducible triadic $3$-cocycle is exhibited
there.  We report this as a failure of the finite model, not as a success in
disguise.

\subsection{Part II in one paragraph}

Part~II takes seriously the place where a reading is undetermined.  It makes the
typed hole --- the \emph{aperture} --- the primitive: syntactically empty,
semantically a family of fillings, and with open and close operations that are
deliberately not inverse.

On that primitive it builds a finite semantics (a terminal-world carrier, filling
fibres, and an adequacy theorem for threaded logic under explicit hypotheses), an
ordered linear natural-deduction fragment that is sound but provably incomplete for
Boolean semantics, and a proof theory strong enough to be interesting: an ordered
one-hole substitution law with a ledger splice, conditional fresh-cut admissibility,
strong normalization via two decreasing measures, local and global confluence with a
unique normal form per fixed proof, and finally a focused normal-form characterization
with the subformula property that yields a terminating, sound, and complete decision
procedure for derivability.

Part~II closes with the six-port whole-cut kernel, which answers the question ``why
six ports?'' by a counting theorem $n\ge 2m$ under a stated contract, and which
weakens universality to relative initiality rather than claiming it.

\subsection{The discipline the two parts share}

Three rules govern every claim in this paper, and the reader is invited to hold the
paper to them.

\begin{enumerate}
\item \textbf{A quotient is admissible only with a declared residual.}  Every
  quotient result must either be accompanied by an exact reconstruction statement or
  name a residual whose contents are characterized.  The phrase ``the quotient loses
  nothing important'' is not an available move in this paper.
\item \textbf{A calibration is never promoted to a theorem.}  Where a statement is
  established only by an exhaustive check of a declared finite fixture, it is
  labelled a calibration, and the fixture is named.  Where it is established by an
  argument, the argument appears in the paper or in an appendix.
\item \textbf{A property is certified, not named.}  Three readings are triadic only
  if the pairwise crossings are given and the cyclic residual is exhibited.
  Universality is only relative initiality unless a stronger theorem is supplied.
  A boundary-return word is not a periodic orbit.
\end{enumerate}

These rules are the reason Part~I contains as many no-go results as positive ones.
They are also the reason Part~II spends effort on metatheorems: soundness alone is
cheap, and the interesting question is always what fails.

\subsection{Reading guide and dependencies}

A reader interested only in the geometry of readings may read Part~I, taking
\S\ref{sec:observers}--\S\ref{sec:triadic-inference} as the core and
\S\ref{sec:classical}--\S\ref{sec:through} as the calibrations that constrain it.
A reader interested in the syntax of apertures may read Part~II from
\S\ref{sec:apertures}, taking \S\ref{sec:ordered-linear}--\S\ref{sec:confluence-decidability}
as its proof-theoretic content.  The two parts meet at the notion of a filling fibre:
in Part~I a residual is what a quotient cannot see, and in Part~II an aperture is
where a partial reading records that it cannot yet see.

The paper imports, and does not redefine:

\begin{itemize}
\item from Paper~0, the hyperbolic model, the affine cocycle and its relative defect,
  the M\"obius involution $z\mapsto -1/z$, and the cut-and-glue language
  \cite{Yuan2026AEGPaperZero};
\item from Paper~I, the separation of literal syntax, marked history, induced
  operator, and endpoint value, and the projective-versus-affine distinction
  \cite{Yuan2026AEGFoundations};
\item from Paper~IV, the requirement that a quotient be judged against a declared
  continuation class and observation, and the separation of a semantic bound from an
  operational trace \cite{Yuan2026AEGCondensation}.
\end{itemize}

\subsection{Conventions and notation}

Throughout, $\R$ and $\C$ are the real and complex fields, $\F_2$ the field with two
elements, and $\PP^1$ the projective line.  For a finite set $X$ we write
$\cP(X)$ for its powerset and $\mathbf 2=\{\bot,\top\}$ for the Sierpinski lattice
with $\bot<\top$.

An \emph{observer} is written $Q$; its components are introduced in
\S\ref{sec:observers} and are never abbreviated away.  A \emph{characteristic} is
written $\chi$, and $\Fib_{D,Q}(x)$ denotes an observation fibre; the subscript
always names both the domain and the observer, because omitting either is the most
common source of error in this material.

Crossings are written with subscripts naming the two readings being compared, so
that $J_{\cK\cX}$ and $J_{\cX t}$ are different maps and composition order matters.
The orientation expression of the six-state carrier is $\Omega$, and it is a
\emph{derived} object (\S\ref{sec:six-state}) whenever the derivation is available.

In Part~II, a hole is written $h_d$ with $d$ its domain slot; $\Fill_{\cM,\rho}$ is
the filling family in model $\cM$ under assignment $\rho$; $\mathrm{TND}_0$ names the
ordered linear natural-deduction fragment of \S\ref{sec:ordered-linear}.  The
dedicated environments \textsf{Calibration}, \textsf{Proposal}, and
\textsf{Conjecture} carry the epistemic weight that the words suggest, and a
\textsf{Theorem} or \textsf{Proposition} is never used for a finite fixture check.

\subsection{What this paper does not establish}

For emphasis, and because the material is exploratory in places, the following are
\emph{not} claimed anywhere in this paper.

\begin{itemize}
\item That the three named readings are canonically identified.  They are compared
  through declared crossings with declared residuals; a canonical identification is
  an open problem (\S\ref{sec:boundaries}).
\item That arithmetic programs must use complex foundational scalars.  Only that,
  within the tested quadratic family and under a stated bridge hypothesis, one
  property selects $\C$ (\S\ref{sec:scalar-trichotomy}).
\item That the six-state carrier is a spacetime decomposition or a physical claim of
  any kind.  The derivation is finite and combinatorial (\S\ref{sec:six-state}).
\item That any finite calibration demonstrates universality, non-computability, or a
  physical realization.
\item That order can be inferred from an energy, a grade, or a scalar.  The opposite
  is proved in several places (\S\ref{sec:satisfaction-logic},
  \S\ref{sec:residual-catalogue}).
\item That the triadic programme has been completed.  Its decisive missing theorem
  is named in \S\ref{sec:boundaries}.
\end{itemize}
